\documentclass[a4paper,9pt]{extarticle}

%=========================================================
% PACKAGES
%=========================================================

\usepackage[
a4paper,
top=1.8cm,
bottom=1cm,
left=1cm,
right=1cm,
headheight=14pt,
headsep=5mm
]{geometry}

\usepackage[T1]{fontenc}
\usepackage[utf8]{inputenc}
\usepackage[french]{babel}
\usepackage[sfdefault]{FiraSans}

\renewcommand{\familydefault}{\sfdefault}

\usepackage{multicol}
\setlength{\columnsep}{0.7cm}

\usepackage{amsmath,amssymb}
\usepackage{titlesec}
\usepackage{enumitem}

\usepackage{tcolorbox}
\tcbuselibrary{skins,breakable}

\usepackage{xcolor}
\usepackage{graphicx}
\usepackage{float}
\usepackage{tikz}

\usepackage{fancyhdr}
\usepackage[hidelinks]{hyperref}


%=========================================================
% PAGE SETTINGS
%=========================================================

\setlength{\parindent}{0pt}
\setlength{\parskip}{1.5pt}

\setlength{\abovedisplayskip}{4pt}
\setlength{\belowdisplayskip}{4pt}

\setlist[itemize]{
leftmargin=*,
itemsep=1pt,
topsep=1pt
}


%=========================================================
% COLORS
%=========================================================

\definecolor{Blue}{HTML}{E3F0FF}
\definecolor{Green}{HTML}{E3F7EA}
\definecolor{Pink}{HTML}{FFE5ED}
\definecolor{Orange}{HTML}{FFF0DC}
\definecolor{Purple}{HTML}{EEE7FF}
\definecolor{Yellow}{HTML}{FFFADB}

\definecolor{BlueTitle}{HTML}{285B8F}
\definecolor{GreenTitle}{HTML}{287052}
\definecolor{PinkTitle}{HTML}{A84063}
\definecolor{OrangeTitle}{HTML}{A96500}
\definecolor{PurpleTitle}{HTML}{6245A5}


%=========================================================
% HEADER
%=========================================================

\pagestyle{fancy}

\fancyhf{}

\fancyhead[L]{\small\textbf{Computer Security}}
\fancyhead[C]{\small Definitions}
\fancyhead[R]{\small Course Notes}
\fancyhead[L]{\small\textbf{Computer security and privacy}}
\fancyhead[C]{\small Introduction  course - definitions}
\fancyhead[R]{\small Week 01 --- 07-13/09/2026}

\renewcommand{\headrulewidth}{0.4pt}


%=========================================================
% TITLES
%=========================================================

\titleformat{\section}
{\large\bfseries\color{BlueTitle}}
{}{0pt}{}

\titlespacing{\section}
{0pt}{6pt}{3pt}


% Separator

\newcommand{\separator}{
\vspace{2pt}
\hrule
\vspace{3pt}
}


%=========================================================
% BOXES
%=========================================================

% DEFINITION - blue

\newtcolorbox{definition}{
enhanced,
breakable,
colback=Blue,
colframe=BlueTitle,
boxrule=0.7pt,
arc=0mm,
title=\textbf{Definition},
fonttitle=\bfseries
}


% IMPORTANT - pink

\newtcolorbox{important}{
enhanced,
breakable,
colback=Pink,
colframe=PinkTitle,
boxrule=0.7pt,
arc=0mm,
title=\textbf{Key idea},
fonttitle=\bfseries
}


% FORMULAS - green

\newtcolorbox{formula}{
enhanced,
breakable,
colback=Green,
colframe=GreenTitle,
boxrule=0.7pt,
arc=0mm,
title=\textbf{Key formulas},
fonttitle=\bfseries
}


% WARNING - orange

\newtcolorbox{warning}{
enhanced,
breakable,
colback=Orange,
colframe=OrangeTitle,
boxrule=0.7pt,
arc=0mm,
title=\textbf{Warning},
fonttitle=\bfseries
}


% EXAMPLE - purple

\newtcolorbox{example}{
enhanced,
breakable,
colback=Purple,
colframe=PurpleTitle,
boxrule=0.7pt,
arc=0mm,
title=\textbf{Example},
fonttitle=\bfseries
}


% SUMMARY - yellow

\newtcolorbox{summary}{
enhanced,
breakable,
colback=Yellow,
colframe=OrangeTitle,
boxrule=0.7pt,
arc=0mm,
title=\textbf{Summary},
fonttitle=\bfseries
}


%=========================================================
% DOCUMENT
%=========================================================

\begin{document}

\thispagestyle{fancy}

\begin{multicols}{2}


%=========================================================
\section{Computer Security}
%=========================================================

\begin{definition}

\textbf{Computer Security}

\vspace{2pt}

Properties of a computer system must hold in the presence of a
resourced strategic adversary.

\end{definition}


\separator


\textbf{The three main security properties}


\begin{definition}

\textbf{Confidentiality}

\vspace{2pt}

Prevention of unauthorized disclosure of information.

\end{definition}


\begin{example}

The adversary should not be able to read my bank statement.

\end{example}


\begin{definition}

\textbf{Integrity}

\vspace{2pt}

Prevention of unauthorized modification of information.

\end{definition}


\begin{example}

The adversary should not be able to change my bank balance.

\end{example}


\begin{definition}

\textbf{Availability}

\vspace{2pt}

Prevention of unauthorized denial of service.

\end{definition}


\begin{example}

The adversary should not prevent me from accessing my bank account.

\end{example}


%=========================================================
\section{Security Policy}
%=========================================================

\begin{definition}

\textbf{Security Policy}

\vspace{2pt}

A high-level description of the security properties that must hold
in the system in relation to assets and principals.

\end{definition}


\separator


\begin{definition}

\textbf{Assets (Objects)}

\vspace{2pt}

Anything with value that needs to be protected.

\end{definition}


\begin{example}

Examples of assets:

\begin{itemize}

\item data;

\item files;

\item memory;

\item sensitive information.

\end{itemize}

\end{example}


\begin{definition}

\textbf{Principals (Subjects)}

\vspace{2pt}

People, computer programs, services, etc.

Principals may not contain the adversary.

\end{definition}


%=========================================================
\section{Threat Model}
%=========================================================

\begin{definition}

\textbf{Threat Model}

\vspace{2pt}

Describes the resources available to the adversary and the
adversary's capabilities.

\end{definition}


\textbf{Typical adversary capabilities}

\begin{itemize}

\item observe;

\item influence;

\item corrupt;

\item exploit available resources.

\end{itemize}


\begin{definition}

\textbf{Adversary}

\vspace{2pt}

The adversary is a malicious entity aiming at breaching the
security policy.

The adversary is strategic: the adversary chooses the optimal way
to use their resources to mount an attack that violates the
security properties.

\end{definition}


%=========================================================
\section{Vulnerabilities and Threats}
%=========================================================

\begin{definition}

\textbf{Vulnerability}

\vspace{2pt}

A specific weakness that could be exploited by adversaries with
interest in different assets.

\end{definition}


\begin{example}

Examples:

\begin{itemize}

\item The banking API is not protected.

\item The password appears in plain text on my screen.

\end{itemize}

\end{example}


\separator


\begin{definition}

\textbf{Threat}

\vspace{2pt}

The feared event or the goal of the adversary that we do not want
to materialize.

\end{definition}


\begin{example}

Examples:

\begin{itemize}

\item A hacker wants to retrieve money by breaking into the bank's
system.

\item A student wants to learn my password by looking over my
shoulder.

\end{itemize}

\end{example}


\begin{definition}

\textbf{Harm}

\vspace{2pt}

The bad thing that happens when the threat materializes.

\end{definition}


\begin{example}

Examples:

\begin{itemize}

\item The adversary steals money.

\item The adversary blocks access to the bank.

\item The adversary learns my password.

\item The adversary reads a message.

\end{itemize}

\end{example}


%=========================================================
\section{Relationship Between Concepts}
%=========================================================

\begin{important}

If the security properties hold within the threat model, it means
that:

\begin{itemize}

\item there are no vulnerabilities that can be exploited;

\item threats cannot materialize;

\item no harm happens.

\end{itemize}

\end{important}


\separator


%=========================================================
\section{Security Mechanisms}
%=========================================================

\begin{definition}

\textbf{Security Mechanism}

\vspace{2pt}

A technical mechanism used to ensure that the security policy is
not violated by an adversary within the threat model.

\end{definition}


\begin{important}

Security mechanisms protect the system against adversaries whose
resources and capabilities are described by the threat model.

\end{important}


%=========================================================
\section{When is a System Secure?}
%=========================================================

\begin{definition}

\textbf{Secure System}

\vspace{2pt}

A system is secure if an adversary constrained by a specific threat
model cannot violate the security policy.

\end{definition}


\begin{summary}

A secure system depends on:

\begin{itemize}

\item a clearly defined security policy;

\item identified assets and principals;

\item a realistic threat model;

\item effective security mechanisms.

\end{itemize}

\end{summary}


%=========================================================
\section{Security Argument}
%=========================================================

\begin{definition}

\textbf{Security Argument}

\vspace{2pt}

A rigorous argument that the security mechanisms in place are
indeed effective in maintaining the security policy.

The argument can be verbal or mathematical and is subject to the
assumptions of the threat model.

\end{definition}



\end{multicols}

\end{document}